% mazzi: 03
\chapter{I presupposti scientifici della valutazione funzionale}

\textbf{Presupposti scientifici}: gli aspetti su cui si fonda la valutazione funzionale, e quindi
il punto di partenza per sviluppare le \textbf{metodologie} con cui indagare le caratteristiche
fisiche e motorie ritenute importanti.

% ---------------------------------------------------------------------------
\section{Obiettivi della valutazione funzionale}

\begin{enumerate}
  \item \textbf{indagine sulle qualità fisiche} --- organiche, neuromuscolari, metaboliche --- che
    influenzano la prestazione;
  \item \textbf{controllo ed ottimizzazione dell'allenamento}, in termini di volume e intensità:
    \begin{itemize}
      \item carichi di lavoro eccessivi per un soggetto possono risultare ottimali per un altro
        $\rightarrow$ occorre passare attraverso test diversi per indirizzare il lavoro;
      \item il lavoro va indirizzato in modo \textbf{individuale}, secondo l'esigenza del singolo
        giocatore, e \textbf{specifico}, secondo la disciplina e il ruolo ricoperto;
    \end{itemize}
  \item \textbf{diagnosi funzionale}, cioè controllo degli \textbf{adattamenti biologici}:
    \begin{itemize}
      \item l'allenamento va inteso come \textbf{programmazione di stimoli meccanici};
      \item \textbf{variazione ambientale} $\rightarrow$ una struttura biologica viene sottoposta a
        stimoli meccanici, di cui occorre poi controllare le \textbf{risposte adattive};
    \end{itemize}
  \item \textbf{ricerca ed identificazione del talento}:
    \begin{itemize}
      \item il docente osserva però che sarebbe più corretto parlare di \textbf{inclinazioni
        attitudinali} verso alcuni sport o alcuni ruoli;
      \item il talento non ha bisogno di essere valutato: risalta subito, si fa notare;
      \item i test servono semmai, per tutti coloro che non sono talenti, a individuare le
        inclinazioni verso determinate discipline o determinati ruoli;
    \end{itemize}
  \item \textbf{ricerca scientifica}: unico modo per sviluppare \textbf{protocolli di ricerca} sul
    \textbf{modello di prestazione} e sui \textbf{metodi di allenamento}, per capire se abbiano o
    meno validità.
\end{enumerate}

% ---------------------------------------------------------------------------
\section{Il modello funzionale della prestazione}

Per arrivare al modello funzionale della prestazione esistono due strade.

\begin{enumerate}
  \item \textbf{valutazione dei parametri fisiologici} misurabili durante la prestazione:
    \begin{itemize}
      \item è la via più complicata, benché la tecnologia stia aiutando sempre di più: oggi si
        dispone di strumenti \textbf{miniaturizzati}, tali da non influenzare la prestazione;
      \item resta difficile negli \textbf{sport da contatto} come il calcio, o comunque in partita
        per alcuni parametri;
    \end{itemize}
  \item \textbf{valutazione delle caratteristiche fisiologiche} dell'atleta che pratica quella
    disciplina sportiva:
    \begin{itemize}
      \item si fonda sul fatto che la disciplina praticata condiziona in maniera specifica le
        capacità fisico-organiche del giocatore;
      \item è probabilmente l'unica strada praticabile negli sport da contatto, quindi anche nel
        calcio.
    \end{itemize}
\end{enumerate}

Ne consegue l'\textbf{identificazione del modello di prestazione} attraverso la valutazione delle
caratteristiche \textbf{metaboliche} e \textbf{neuromuscolari} tipiche dell'atleta.

% ---------------------------------------------------------------------------
\section{Classificazione delle discipline sportive}

Attraverso le \textbf{richieste organico-funzionali in gara} si identifica il tipo e la maggior
incidenza di ciascuna di esse nel determinare la prestazione.

\begin{itemize}
  \item discipline a impegno \textbf{prevalentemente aerobico};
  \item discipline a impegno \textbf{aerobico-anaerobico massivo};
  \item discipline a impegno \textbf{prevalentemente anaerobico};
  \item discipline a impegno \textbf{aerobico-anaerobico alternato};
  \item discipline di \textbf{potenza};
  \item discipline di \textbf{destrezza}.
\end{itemize}

\rubrica{Limite di queste classificazioni}
La letteratura scientifica è piena di classificazioni di questo tipo, ma far rientrare il
\textbf{calcio} in una sola di esse è \textbf{improponibile}: a eccezione forse della prima
(impegno prevalentemente aerobico), il calcio si ritrova sostanzialmente in tutte le altre.

% ---------------------------------------------------------------------------
\section{Definizione di test}

\textbf{Test}: unità essenziale della valutazione funzionale motoria.

\begin{itemize}
  \item la richiesta di un \textbf{compito motorio specifico}, in relazione alla sua modalità di
    esecuzione, coinvolgerà capacità fisiche specifiche;
  \item la valutazione si realizza attraverso la misurazione diretta o indiretta di parametri
    fisici, biomeccanici e/o biologici connessi:
    \begin{itemize}
      \item al compito motorio richiesto;
      \item alla sua modalità di esecuzione;
      \item allo strumento di valutazione utilizzato;
    \end{itemize}
  \item le procedure di esecuzione e di misurazione ne definiscono il \textbf{protocollo}.
\end{itemize}

\rubrica{Esempio: la forza esplosiva degli arti inferiori}
Volendo misurare la forza esplosiva degli arti inferiori, l'unità essenziale --- il test ---
potrebbe essere il \textbf{salto verticale}, realizzabile in due modalità differenti:
\begin{itemize}
  \item con le \textbf{mani ai fianchi}, per evitare l'influenza della componente coordinativa;
  \item con le \textbf{braccia libere}.
\end{itemize}

\rubrica{Che cosa deve contenere il protocollo}
Nel protocollo devono rientrare entrambi gli elementi seguenti, perché è la metodica a rendere i
dati confrontabili tra loro:
\begin{enumerate}
  \item la \textbf{modalità di esecuzione} del gesto;
  \item il \textbf{tipo di strumentazione} utilizzata: usare una \textbf{piattaforma di forza} o un
    \textbf{accelerometro} cambia il protocollo, e possono cambiare anche i risultati.
\end{enumerate}

Ciò vale sia per il confronto \textbf{longitudinale} (pre-post del singolo giocatore) sia per
quello \textbf{trasversale} (tra gruppi di giocatori).

\subsection{Caratteristiche dei test}

I test si dividono in funzione del luogo in cui vengono effettuati, in \textbf{test da
laboratorio} e \textbf{test da campo}; ciascuno dei due si articola a sua volta in generali e
specifici, e tutte le categorie confluiscono nella valutazione del \textbf{sistema metabolico} e
del \textbf{sistema neuromuscolare}. Le associazioni prevalenti sono laboratorio $\rightarrow$
generali e campo $\rightarrow$ specifici.

\begin{center}
\begin{forest} classalbero
[Test
  [Da laboratorio [Generali][Specifici]]
  [Da campo [Generali][Specifici]]
]
\end{forest}
\end{center}

\rubrica{Test da laboratorio}
\begin{itemize}
  \item molto in uso negli anni precedenti, quando le strumentazioni erano ingombranti e
    difficilmente trasportabili;
  \item la valutazione era perciò legata alla logistica dello spazio del laboratorio;
  \item perdevano di specificità: davano informazioni importanti dal punto di vista generale, ma
    poco specifiche;
  \item oggi si stanno occupando di situazioni particolari, riguardanti più l'aspetto
    \textbf{clinico} che quello \textbf{prestativo}.
\end{itemize}

\rubrica{Test da campo}
\begin{itemize}
  \item risultano molto più specifici, perché permettono di utilizzare il gesto più vicino a
    quello della prestazione;
  \item lo strumento un tempo disponibile sul campo era solamente il \textbf{cronometro}; oggi si
    dispone di strumentazioni scientifiche valide, capaci di dare informazioni specifiche sia sul
    sistema neuromuscolare sia sul sistema metabolico.
\end{itemize}

% ---------------------------------------------------------------------------
\section{I presupposti scientifici di un test}

Un test deve possedere quattro requisiti:
\begin{enumerate}
  \item \textbf{validità}: capacità di misurare caratteristiche specifiche; un test è valido se
    misura \textbf{quella} grandezza biologica o capacità fisica che ci si è prefissati di
    valutare;
  \item \textbf{riproducibilità}: attendibilità di risultati simili in prove successive;
  \item \textbf{obiettività}: non deve essere influenzato dall'\textbf{operatore} e/o da
    \textbf{fattori ambientali}; se questi influissero, i risultati non riguarderebbero più
    l'atleta ma altri eventi $\rightarrow$ conclusioni errate;
  \item \textbf{specificità}: solo attraverso di essa si può programmare un allenamento adeguato e
    indagare la capacità fisica che davvero si avvicina al \textbf{modello di prestazione}.
\end{enumerate}

% ---------------------------------------------------------------------------
\section{La validità}

\subsection{Validità sostanziale o di costrutto}

\textbf{Validità sostanziale} (o \textbf{di costrutto}): validità di primaria importanza, che
identifica l'abilità di un test a rappresentare i presupposti teorici che spiegano alcuni aspetti
derivanti dalla conoscenza e dalle osservazioni.

\begin{itemize}
  \item rappresenta il \textbf{limite della misura} di un test nella valutazione delle
    caratteristiche fisiche per le quali è stato progettato;
  \item p.\,es.\ per misurare la forza esplosiva degli arti inferiori occorre un test in grado di
    restituire \textbf{quel} risultato in maniera specifica, senza compromessi.
\end{itemize}

\subsection{Validità formale o apparente}

\textbf{Validità formale} (o \textbf{apparente}): validità di secondaria importanza, che identifica
il \textbf{grado di apparenza esterna} che un test possiede nel misurare le caratteristiche fisiche
e biologiche che si propone di misurare.

\begin{itemize}
  \item è un presupposto di tipo \textbf{qualitativo}: una caratteristica informale e non
    quantitativa;
  \item il soggetto che esegue il test lo riconosce come vicino alla prestazione $\rightarrow$ più
    motivato a eseguire il test o la batteria di test;
  \item la \textbf{motivazione} è essa stessa un aspetto fondamentale della validità del dato: deve
    essere altissima, altrimenti il dato non ha alcun valore;
  \item p.\,es.\ misurare un parametro metabolico di un calciatore usando una \textit{bike} ha una
    validità apparente bassissima: il giocatore capisce da sé che quel tipo di valutazione non ha
    nesso con la sua prestazione $\rightarrow$ la motivazione ne risente.
\end{itemize}

\subsection{Validità di contenuto}

\textbf{Validità di contenuto}: identifica l'abilità di un test e/o di una batteria di test di
misurare tutte le abilità motorie necessarie per una determinata disciplina sportiva.

\begin{itemize}
  \item un \textbf{solo test} non basta a valutare tutte le capacità motorie necessarie allo
    sviluppo della prestazione $\rightarrow$ occorre una \textbf{batteria} di test, che dia un
    quadro più completo di come il giocatore sviluppa la propria prestazione (forza esplosiva,
    \textit{stiffness}, capacità anaerobiche, \dots);
  \item l'importanza data alle abilità motorie indagate deve essere correlata alle richieste del
    compito della disciplina sportiva praticata.
\end{itemize}

\subsection{Validità del criterio di riferimento}

\textbf{Validità del criterio di riferimento}: mette in correlazione la misura di un test con
altri test che misurano la stessa abilità. Se ne distinguono quattro tipi.

\begin{enumerate}
  \item \textbf{validità concorrente}: grado della misura di un test e della strumentazione
    utilizzata, associato (stimato statisticamente) a quella di un altro test e della relativa
    strumentazione, già accettati (ritenuti validi) per la misura della stessa abilità;
    \begin{itemize}
      \item si ha quando due test si avvicinano l'uno all'altro sia per la \textbf{capacità motoria
        indagata} sia per la \textbf{strumentazione utilizzata};
    \end{itemize}
  \item \textbf{validità convergente}: \textbf{alta correlazione} tra le misure di un test e quelle
    di un test riconosciuto come \textit{gold standard};
    \begin{itemize}
      \item p.\,es.\ il \textbf{salto in lungo}, considerato espressione della forza esplosiva, ha
        validità convergente con il \textbf{salto verticale};
    \end{itemize}
  \item \textbf{validità predittiva}: limite della misura di un test in relazione alla misura
    ottenuta nello stesso da \textbf{atleti di alto livello} praticanti la stessa disciplina
    sportiva, cioè da coloro che rappresentano il \textbf{modello di prestazione} $\rightarrow$ la
    misura può essere predittiva della prestazione futura;
  \item \textbf{validità discriminante}: abilità di un test di distinguersi tra due differenti
    costrutti;
    \begin{itemize}
      \item si evidenzia con una scarsa correlazione tra i risultati del test rispetto a quelli
        di un altro costrutto;
      \item buona validità discriminante in una batteria $\rightarrow$ evita inutili spese di
        tempo, energie e risorse nella somministrazione di test altamente correlati tra loro;
      \item i test di una batteria devono misurare in modo indipendente le abilità motorie
        componenti la prestazione;
      \item p.\,es.\ \textbf{salto in lungo} e \textbf{salto in alto} non hanno motivo di coesistere
        nella stessa batteria: indagano la medesima capacità $\rightarrow$ sovrapposizione inutile
        di dati, tempo ed energie.
    \end{itemize}
\end{enumerate}

% ---------------------------------------------------------------------------
\section{La riproducibilità}

\textbf{Riproducibilità}: misura del \textbf{grado di consistenza} o di riproducibilità di un test.

\begin{itemize}
  \item se un atleta possiede un'abilità indagata dal test, i risultati non devono cambiare in
    due misure successive;
  \item il test è riproducibile solo se si ottengono le stesse misure con \textbf{scarsissima
    variabilità} in due momenti successivi di somministrazione;
  \item una variabilità fra misure eseguite a distanza di un giorno che non sia imputabile ad
    alcun fattore esterno, ma solo all'esecuzione del test, rende il test non riproducibile;
  \item un test deve essere \textbf{valido e riproducibile}, perché un'elevata variabilità non
    avrebbe senso.
\end{itemize}

\rubrica{Riproducibilità test-retest}
\textbf{Test-retest}: amministrare il test due volte allo stesso gruppo di atleti ed effettuare
la correlazione statistica tra le due misure.

\subsection{Validità e riproducibilità a confronto}

I due requisiti sono \textbf{indipendenti}, e se ne danno quattro combinazioni.

\begin{enumerate}
  \item \textbf{valido ma non riproducibile}: la \textbf{media} delle misure ripetute si avvicina
    al valore reale, ma le singole misure presentano una variabilità molto alta;
  \item \textbf{riproducibile ma non valido}: le misure ripetute sono tutte vicine o uguali tra
    loro --- scarsissima variabilità --- ma lontane dal valore reale atteso;
  \item \textbf{né valido né riproducibile}: le misure sono distanti tra loro e distanti dal valore
    di riferimento;
  \item \textbf{valido e riproducibile}: non solo c'è scarsa variabilità tra le misure prese
    singolarmente, ma tutte si avvicinano al valore reale dell'abilità indagata.
\end{enumerate}

% ---------------------------------------------------------------------------
\section{La specificità}

\textbf{Specificità di una valutazione}: caratteristica di un test di basarsi, \textbf{laddove
possibile}, su condizioni simili a quelle ambientali e fisiche della competizione, cioè sulle
dinamiche dell'azione e del gesto sportivo.

\begin{itemize}
  \item il test è \textbf{specifico} se riproduce e soddisfa le condizioni della prestazione in
    termini di \textbf{coerenza biomeccanica}, coordinativa e metabolica.
\end{itemize}

\rubrica{Test funzionali specifici}
I \textbf{test funzionali specifici} mirano alla diagnosi di qualità fisiologiche, neuromuscolari
e biomeccaniche tipiche di una disciplina sportiva.

\rubrica{Esempio: le catene cinetiche nel gesto del calciatore}
Trattandosi di movimenti alquanto complessi, saper organizzare test specifici permette di
distinguere le azioni diverse dei due arti nello stesso gesto:
\begin{itemize}
  \item \textbf{arto calciante} $\rightarrow$ catena cinetica \textbf{aperta};
  \item \textbf{arto di appoggio} $\rightarrow$ catena cinetica \textbf{chiusa}, con la funzione
    completamente diversa di mantenere l'equilibrio del corpo.
\end{itemize}

La distinzione risulta interessante non tanto ai fini della prestazione, quanto nella \textbf{fase
di recupero post-infortunio}.
